\section{Marco Teórico}

\subsection{Teoría sobre la creación de un inicio de sesión}
1.1
La autenticación de usuarios es el mecanismo fundamental que permite a un sistema operativo validar la identidad de los sujetos que solicitan acceso a sus recursos. De acuerdo con Silberschatz, Galvin y Gagne en Fundamentos de Sistemas Operativos, es indispensable establecer una distinción clara entre la identificación y la autenticación. La identificación ocurre cuando un usuario declara su identidad mediante una credencial o nombre de cuenta. Por su parte, la autenticación comprende la comprobación explícita de dicha declaración. Las técnicas de autenticación se dividen fundamentalmente en tres categorías: información que el usuario conoce, objetos que el usuario posee e indicadores biométricos que identifican sus características físicas o conductuales.

En el contexto de la administración de accesos en entornos tipo UNIX, Francisco Manuel Márquez García describe en Unix Programación Avanzada la secuencia técnica de llamadas al sistema encargadas del flujo de inicio de sesión. La arquitectura de estos sistemas delega la gestión de las líneas de terminal al proceso inicializador del sistema, el cual invoca rutinas especializadas para configurar la línea de comunicación y solicitar el identificador del usuario. Posteriormente, se transfiere el control al ejecutable de inicio de sesión, el cual desactiva el eco en la terminal, solicita la clave de acceso y consulta los archivos de seguridad del sistema antes de inicializar el entorno mediante el establecimiento de identificadores de usuario y la ejecución del intérprete de comandos.

La arquitectura de protección de credenciales en el sistema operativo exige una separación rigurosa en el almacenamiento de datos. Como exponen Andrew S. Tanenbaum y Herbert Bos en Sistemas Operativos Modernos, las contraseñas nunca deben almacenarse en texto plano. Históricamente, la información técnica de los usuarios se concentraba en el archivo del sistema de acceso público; sin embargo, los sistemas modernos implementan mecanismos de protección que restringen los permisos de lectura únicamente al usuario con privilegios de administrador, albergando los resúmenes cifrados de las contraseñas junto con los parámetros de expiración de las cuentas.

1.1.1. Identificación y autenticación de usuarios
El proceso de identificación constituye el primer paso dentro del control de acceso en un sistema operativo multiusuario. En esta etapa, el usuario proporciona un identificador único, como un nombre de usuario o un número de cuenta, para declararse ante el sistema. No obstante, la sola declaración de la identidad es insuficiente para otorgar privilegios sobre los recursos del sistema operativo, por lo que se requiere un procedimiento riguroso de autenticación. Silberschatz, Galvin y Gagne destacan que la autenticación verifica la validez de la afirmación presentada. La forma más extendida de autenticación se basa en secretos compartidos, como las contraseñas numéricas o alfanuméricas, donde el usuario demuestra su conocimiento sobre un dato privado previamente registrado. En entornos que requieren mayor grado de seguridad, los sistemas operativos integran múltiples factores combinados para reducir las vulnerabilidades derivadas del robo o la adivinación de claves de acceso.

1.1.2. El flujo de inicio de sesión en sistemas UNIX
En los sistemas operativos basados en la filosofía UNIX, el inicio de sesión es administrado mediante una jerarquía de procesos encadenados. Márquez García explica que el proceso raíz o inicializador del sistema genera procesos dedicados a supervisar cada terminal o línea de entrada disponible. Cuando un usuario interactúa con la terminal, la rutina getty establece la velocidad de transmisión, los caracteres de control de la línea y despliega el mensaje de bienvenida solicitando el nombre de usuario. Tras la recepción del nombre de usuario, dicho proceso sustituye su imagen en memoria mediante la llamada al sistema execve para invocar al programa login. El ejecutable login se encarga de cambiar la configuración del controlador de la terminal para deshabilitar el eco del teclado, solicita la contraseña, aplica el algoritmo de verificación y, en caso de autenticación exitosa, asigna los identificadores de usuario y grupo correspondientes mediante llamadas al sistema como setuid y setgid, finalizando con la ejecución del shell configurado para el usuario.

1.1.3. Archivos de control y protección de credenciales (/etc/passwd y /etc/shadow)
La estructura de almacenamiento de credenciales en sistemas UNIX evolucionó debido a las limitaciones de seguridad del diseño original. Tanenbaum y Bos detallan que tradicionalmente el archivo /etc/passwd guardaba todos los datos de la cuenta de usuario, incluyendo el nombre, la contraseña cifrada, el identificador de usuario, el identificador de grupo predeterminado, el directorio de trabajo y el intérprete de comandos asignado. Dado que numerosos programas legítimos del sistema necesitan consultar la información de los usuarios en /etc/passwd, dicho archivo debe conservar permisos de lectura global para todos los procesos del sistema. Esta visibilidad pública permitía que usuarios malintencionados copiaran el archivo y aplicaran técnicas de fuerza bruta o diccionarios fuera de línea sobre las contraseñas cifradas. Para erradicar este riesgo, se introdujo la técnica del archivo en sombra o /etc/shadow, el cual separa las contraseñas cifradas y las traslada a un archivo protegido accesible únicamente por el usuario administrador. En /etc/passwd el campo de la contraseña se sustituye por una marca fija, mientras que /etc/shadow resguarda los datos sensibles, la sal de cifrado y los parámetros de expiración de la clave.


\subsection{Algoritmos de cifrado}
El almacenamiento de contraseñas en texto plano representa una vulnerabilidad crítica. Para mitigar este riesgo, los sistemas implementan algoritmos de cifrado o funciones de un solo sentido (hashing). En lenguaje C, es común utilizar funciones matemáticas o bibliotecas nativas como \texttt{crypt()} (basada en DES, MD5 o SHA) que transforman la contraseña original en una cadena de caracteres ilegible. Durante el inicio de sesión, el sistema cifra la contraseña ingresada y compara el resultado con el hash guardado; si ambos coinciden, se otorga el acceso sin que el sistema necesite conocer la clave original.

\subsection{Biblioteca ncurses}
\texttt{ncurses} (new curses) es una biblioteca de programación que proporciona una API para construir interfaces de usuario basadas en texto de manera independiente de la terminal. Permite controlar la posición del cursor, crear ventanas, utilizar colores y leer el teclado carácter por carácter. Funciones como \texttt{noecho()} son esenciales en los inicios de sesión, ya que desactivan el eco visual de la terminal, permitiendo que el usuario escriba su contraseña sin que los caracteres sean visibles en la pantalla.

\newpage
